\chapter{The Breathing-Space of NEP}

War communism had been made up of two major elements: on the one hand, a concentration of economic authority and power, including centralized control and management, the substitution of large for small units of production and some measure of unified planning; on the other hand, a flight from commercial and monetary forms of distribution, and the introduction of the supply of basic goods and services free or at fixed prices, rationing, payments in kind, and production for direct use rather than for a hypothetical market. Between these two elements, however, a fairly clear distinction could be drawn. The processes of concentration and centralization, though they flourished exceedingly in the forcing-house of war communism, were a continuation of processes already set in motion during the first period of the revolution, and indeed during the European war. Here war communism was building on a foundation of what had gone before, and many of its achievements stood the test; only in their detailed application were its policies afterwards subject to rejection and reversal. The second element of war communism, the substitution of a ``natural'' for a ``market'' economy, had no such foundations. Far from developing logically out of the policies of the initial period of the revolution, it was a direct abandonment of those policies---an unprepared plunge into the unknown. These aspects of war communism were decisively rejected by NEP; and it was these aspects which most of all discredited it in the eyes of its critics. 

Between the two major elements of war communism there was, moreover, a further distinction. The policies of concentration and centralization were applied almost exclusively in industry; attempts to transfer them to agriculture met with no success. It was here that the revolution had the main social basis of its support, and that the Russian economy showed some of the features of a developed capitalism. The policies of the flight from money and the substitution of a ``natural'' economy arose, not from any preconceived plan, but from inability to solve the problems of a backward peasant agriculture which occupied more than 80 per cent of the population. They were an expression of the fundamental difficulty of attempting to run in double harness the anti-feudal revolution of a peasantry with petty-bourgeois aspirations and the anti-bourgeois, anti-capitalist revolution of a factory proletariat, and of coping with the conflict between town and country inherent in the attempt. These were the incompatibilities which eventually brought the revolt against war communism and destroyed it. 
% 注：去除了 "-7^"、"-7""、"^\j"、"Ition"、"!a"、"jbilities" 等乱码。

By the autumn of 1920, when the fighting was over, the whole economy was grinding to a halt. Nothing in the theory or practice of war communism offered any clue how to re-start processes of production and exchange which had come to a standstill. The nodal point, as always in the Russian economy, was grain. The policy of requisitions, which had worked after a fashion during the civil war, was bankrupt. The peasant retreated into a subsistence economy, and had no incentive to produce surpluses which would be seized by the authorities. Widespread peasant disorders occurred in central Russia during the winter of 1920-1921. Gangs of demobilized soldiers roamed the countryside in search of food, and lived by banditry. It was imperative, if the rest of the country was not to starve, to provide the peasant with the incentives which were denied to him under a system of requisition. Nor was all well within the party. A dissentient group calling itself a ``Workers' Opposition'' was formed under the leadership of Shlyapnikov, a former metal worker who had been People's Commissar for Labour in the first Soviet Government, and Alexandra Kollontai, who enjoyed some prestige in the early days of the revolution. Its programme was directed mainly against the proliferation of economic and political controls and the growing power of the party and state machine; it claimed to uphold the purity of the original ideals of the revolution, and recalled the opposition of 1918 to the surrender of Brest-Litovsk. The leadership of the group was not very impressive. But it enjoyed wide sympathy and support in the party ranks. 
% 注：去除了 ". to a standstill" 前的杂乱标点，清理了 "---r"、"tjjof"、"^sympathy"。

A change of front was now urgently necessary. The essence of the new policy worked out during the winter of 1920-1921 was to permit the peasant, after the delivery of a fixed proportion of his output to state organs (a ``tax in kind''), to sell the rest on the market. To make this possible, encouragement must be given to industry, especially small artisan industry, to produce the goods which the peasant would want to buy---a reversal of the emphasis under war communism on large-scale heavy industry. Private trade must be allowed to revive; here much reliance was placed on the cooperatives---one of the few pre-revolutionary institutions to retain some degree of vitality and popularity. Finally, all this implied---though the point was not grasped till somewhat later---a halt to the headlong fall of the ruble and the establishment of a stable currency. The package known as the New Economic Policy (NEP), with particular emphasis on concessions to the peasant, was approved by the central committee for presentation by Lenin to the historic tenth party congress in March 1921. 
% 注：去除了 "cwas", "jlproportion", "|fto", "Vto", "da", "llhalt", "//of" 等乱码。

On the eve of the congress its proceedings were overshadowed by a sinister and ominous disaster. The sailors of the Red fleet based on the fortress of Kronstadt rose in revolt, demanding concessions for workers and peasants and the free election of Soviets. The mutiny had no direct association with the Workers' Opposition, but reflected the same deep feelings of discontent with the trend of party policy. Such leadership as there was appears to have been anarchist; the suspicion of the Bolsheviks that it had been planned or inspired by White emigres was unfounded, though they afterwards made much capital out of it. Parleys and calls to surrender were fruitless. On March 17, while the congress was debating Lenin's proposals, units of the Red Army advanced on the fortress across the ice. After a bloody battle, fought on both sides with great tenacity, the rebels were overpowered and the fortress seized. But this massive revolt of men hitherto honoured as heroes of the revolution was a staggering blow to the prestige and confidence of the party. It may well have increased the readiness of the congress to accept the New Economic Policy, as well as proposals to tighten party discipline and provide stronger safeguards against dissent, within and without. 
% 注：去除了 "iWorkers'", "/bf discontent", "hfortress", "fi prestige" 等乱码。

When Lenin submitted the resolution embodying the NEP proposals to the congress, the debate was perfunctory. Disenchantment with war communism was general; and the crisis was too acute to brook delay. Doubters were consoled by Lenin's assurance that ``the commanding heights'' of industry would remain in the firm hands of the state, and that the monopoly of foreign trade would be maintained intact. The resolution was accepted, if not with enthusiasm, with good grace and formal unanimity. The sharpest difference of opinion at the congress arose out of the heated controversy on the trade union question which had raged throughout the winter. Trotsky, inspired by the experience of the civil war, and supported after some hesitation by Bukharin, once more propounded his plan for transforming the trade unions into ``production unions'', and making them part of ``the apparatus of the workers' state''. At the opposite end of the spectrum, the Workers' Opposition wished to place the organization and control of production in the hands of the workers, as represented in the trade unions---a quasi-syndicalist view. Manoeuvring between the two embattled factions, Lenin finally succeeded in rallying the centre of the party round a resolution which, however, skirted the crucial issues without solving them. The taint of ``militarization'' was avoided. The trade unions were recognized as ``mass non-party organizations'' which had to be won over. It would be an error to incorporate them in the state machine. Persuasion, not force, was their proper instrument, though ``proletarian compulsion'' was not ruled out. The trade unions had always professed a concern for production; as early as 1920 the trade union central council established a Central Labour Institute for the study of, and training in, methods and techniques designed to improve the productivity of workers. This aspect of their responsibilities was emphasized in the resolution. It was their function to maintain labour discipline and combat absenteeism; but this should be done through ``comradely courts'', not by organs of the state. The resolution was carried by a large majority, but not without some minority votes being cast for two dissentient drafts. 

The bitterness of the controversy shocked the party, and left its mark on the congress. Lenin spoke of the ``fever'' which had shaken the party, and of ``the luxury of discussions'' and ``disputes'' which the party could ill afford. The congress adopted a special resolution bearing the title ``On the Syndicalist and Anarchist Deviation in our Party'', which declared dissemination of the programme of the Workers' Opposition to be incompatible with party membership, as well as a general resolution ``On the Unity of the Party''. This demanded ``the complete abolition of all fractionalism''; disputed issues could be discussed by all members of the party, but the formation of groups with ``platforms'' of their own was banned. Once a decision had been taken, unconditional obedience to it was obligatory. Infringement of this rule could lead to expulsion from the party. A final clause, which was kept secret and published only three years later, laid it down that even members of the party central committee could be expelled on these grounds by a majority of not less than two-thirds of members of the committee. These provisions, designed to ensure loyalty and uniformity of opinion in the party, seemed necessary and reasonable at the time. As Lenin put it, ``during a retreat discipline is a hundred times more necessary''. But the vesting of what was in effect a monopoly of power in the central organization of the party was to have far-reaching consequences. Lenin at the height of the civil war had acclaimed ``the dictatorship of the party'', and maintained that ``the dictatorship of the working class is carried into effect by the party''. The corollary, drawn by the tenth congress, was the concentration of authority in the central organs of the party. The congress conceded to the trade unions a measure of autonomy vis-a-vis the organs of the workers' state. But the role which they were to play was determined by the monopoly of power vested in the party organization. 
% 注：去除了 "jwas banned", "[obedience", "itwo-thirds", "^designed", "[party", "land maintained", "[is carried" 等乱码。将 "vis-d-vis" 修正为 "vis-a-vis"。

The stringent ban on opposition within the party was the product of the crisis which accompanied the introduction of NEP. The same process logically overtook the two Left opposition parties which had survived the revolution; the SRs and the Mensheviks. The dissolution of the Constituent Assembly in January 1918 had proclaimed the determination of the Bolsheviks to exercise supreme power, and laid the foundations of the one-party state. But during the next three years---the period covered by the civil war---mutual relations between the Soviet Government and the two Left parties were ambiguous and fluctuating, and measures taken against them inconclusive. A few weeks after the revolution a group of Left SRs broke away from the main party, and formed a coalition with the Bolsheviks; three Left SRs were appointed People's Commissars. The signature of the Brest-Litovsk treaty in March 1918, which was bitterly denounced both by SRs and by Mensheviks, led to their resignation. The Right SRs now came out openly against the regime, and were held responsible for the disorders in Moscow in the summer of 1918, as well as for the assassination of the German Ambassador and of two leading Bolsheviks in Petrograd, and for the attempt on Lenin's life (see p. 20) above. On June 14, 1918, the Right SRs and the Mensheviks were banned on the ground of their association with ``notorious counter-revolutionaries''. Their newspapers were suppressed from time to time, but often re-appeared under other names; even a Kadet newspaper was published for some months after the revolution. Intermittent harassment, rather than the enforcement of a total ban, reflected ambivalence and hesitation on the part of the authorities. 
% 注：去除了 "nbetween", "/were", "I them", "/Intermittent" 等乱码。

The civil war, which made the plight of the regime more desperate, at first somewhat improved the standing of the two parties. The Mensheviks emphatically, the SRs less consistently, denounced the action of the Whites and of the Allied governments which aided and abetted them, and thus implicitly supported the regime, while continuing to attack its internal policies. The ban on the Mensheviks was withdrawn in November 1918, on the SRs in February 1919; and Menshevik and SR delegates spoke at sessions of the All-Russian Congress of Soviets in 1919 and 1920, though apparently without the right to vote. During the civil war many Mensheviks, and some SRs, joined the Bolshevik party; many more entered the service of the regime and worked in Soviet institutions. The mass following of both parties, persistently harassed by the authorities, began to disintegrate. When the civil war ended, there was no further basis for coalition or compromise. Two thousand Mensheviks, including the whole of the party central committee, were said to have been arrested on the eve of the introduction of the NEP, the extinction of the Menshevik opposition coinciding with the suppression of dissent within the ruling Bolshevik party. Many of those arrested were later released, and the leading Mensheviks allowed to go abroad. But a hard core of SR leaders were put on trial in 1922 for counter-revolutionary activities, and sentenced to death (these sentences were not carried out) or long-term imprisonment. 

The benefits offered by NEP to the peasant, which in any case came too late to affect the sowings for 1921, were retarded by a natural calamity. Severe drought ruined the harvest over a large area, especially in central Russia and in the Volga basin. The famine was more widespread, and worked greater havoc on a much tried and enfeebled population, than the last great Russian famine of 1891. The horrors of the ensuing winter, when millions starved, were partly mitigated by supplies from foreign relief missions, notably the American Relief Administration. Sowings for 1922 were extended. The harvests for that year and for 1923 were excellent, and appeared to herald a revival of Soviet agriculture: small quantities of grain were actually exported. It was remarked that NEP, by re-introducing market processes to the countryside, had reversed the levelling policies of war communism, and encouraged the re-emergence of the rich peasant, or kulak, as the key figure in the rural economy. The poor peasant produced for the subsistence of himself and his family. He consumed what he produced; if he came to the market, it was more often as a buyer than as a seller. The kulak produced for the market and became a small capitalist; this was the essence of NEP. The right to lease land and to employ hired labour, theoretically prohibited since the early days of the regime, was conceded with some formal restrictions in the new agricultural code of 1922. But, so long as the peasants had enough to eat, and provided surpluses sufficient to feed the towns, few even of the most devoted party members were in a hurry to challenge the derogations from the principles and ideals of the revolution which had yielded these fortunate results. If NEP had done little or nothing to help industry or the industrial worker, and less than nothing to promote the cause of a planned economy, these problems could safely be left to the future. 
% 注：去除了 "i go", ", The benefits", "N'W Tw - ty-yt TnnTrn IWny case", "f/harvest", "(mitigated", "iremarked", "jto", "[seller", "Idevoted", "1 and", "'' economy" 等众多识别噪点。

It was at this point that the underlying differences in the party about the character of war communism began to be reflected in differences about the practical implications and consequences of NEP. When, in the crisis atmosphere of March 1921, the substitution of NEP for the more extreme policies of war communism was unanimously accepted as a welcome and necessary relief, these divergences were shelved, but not wholly reconciled. In so far as war communism was thought of, not as an advance on the road to socialism, but as an aberration dictated by military necessities, a forced response to the civil war emergency, NEP was a retracing of steps from a regrettable, though no doubt enforced, digression, and a return to the safer and more cautious path which was being followed before June 1918. In so far as war communism was treated as an over-rash, over-enthusiastic dash forward into the higher reaches of socialism, premature no doubt but otherwise correctly conceived, NEP was a temporary withdrawal from positions which it had proved impossible to hold at the moment, but which would have sooner or later to be regained; and it was in this sense that Lenin, whose position was not altogether consistent, called NEP ``a defeat'' and ``a retreat---for a new attack''. When Lenin at the tenth congress said that NEP was intended ``seriously, and for a long time'' (but added, in reply to a question, that an estimate of 25 years was ``too pessimistic''), he gave hostages both to the view that it was a desirable and necessary correction of the errors of war communism, and to the view that it would itself have in the future to be corrected and superseded. The unspoken premise of the first view was the practical necessity of taking account of a backward peasant economy and peasant mentality; the unspoken premise of the second was the need to build up industry and not further depress the position of the industrial workers who formed the main bulwark of the revolution. These differences, muted for the moment by thankfulness that the acute party crisis of the winter of 1920-1921 had been successfully surmounted, re-appeared in a further economic and party crisis two years later. 
% 注：修复了段尾由于分页造成的断词 "surmounted, . re-appeared"。